%-------------------------
% Resume in Latex
% Author
% License : MIT
%------------------------

%---- Required Packages and Functions ----

\documentclass[a4paper,11pt]{article}
\usepackage{latexsym}
\usepackage{xcolor}
\usepackage{float}
\usepackage{ragged2e}
\usepackage[empty]{fullpage}
\usepackage{wrapfig}
\usepackage{lipsum}
\usepackage{tabularx}
\usepackage{titlesec}
\usepackage{geometry}
\usepackage{marvosym}
\usepackage{verbatim}
\usepackage{enumitem}
\usepackage[hidelinks]{hyperref}
\usepackage{fancyhdr}
\usepackage{fontawesome5}
\usepackage{multicol}
\usepackage{graphicx}
\usepackage{cfr-lm}
\usepackage[T1]{fontenc}
\setlength{\multicolsep}{0pt} 
\pagestyle{fancy}
\fancyhf{} % clear all header and footer fields
\fancyfoot{}
\renewcommand{\headrulewidth}{0pt}
\renewcommand{\footrulewidth}{0pt}
\geometry{left=1.4cm, top=0.8cm, right=1.2cm, bottom=1cm}
\usepackage[most]{tcolorbox}
\tcbset{
	frame code={}
	center title,
	left=0pt,
	right=0pt,
	top=0pt,
	bottom=0pt,
	colback=gray!20,
	colframe=white,
	width=\dimexpr\textwidth\relax,
	enlarge left by=-2mm,
	boxsep=4pt,
	arc=0pt,outer arc=0pt,
}

\urlstyle{same}

\raggedright{}
\setlength{\tabcolsep}{0in}
\setlength{\footskip}{4.5pt}



% Sections formatting
\titleformat{\section}{
  \vspace{-4pt}\scshape\raggedright\large
}{}{0em}{}[\color{black}\titlerule{} \vspace{-7pt}]

%-------------------------
% Custom commands
\newcommand{\resumeItem}[2]{
  \item{
    \textbf{#1}{\hspace{0.5mm}#2 \vspace{-0.5mm}}
  }
}

\newcommand{\resumePOR}[3]{
\vspace{0.5mm}\item
    \begin{tabular*}{0.97\textwidth}[t]{l@{\extracolsep{\fill}}r}
        \textbf{#1}\hspace{0.3mm}#2 & \textit{\small{#3}} 
    \end{tabular*}
    \vspace{-2mm}
}

\newcommand{\resumeSubheading}[4]{
\vspace{0.5mm}\item
    \begin{tabular*}{0.98\textwidth}[t]{l@{\extracolsep{\fill}}r}
        \textbf{#1} & \textit{\footnotesize{#4}} \\
        \textit{\footnotesize{#3}} &  \footnotesize{#2}\\
    \end{tabular*}
    \vspace{-2.4mm}
}

\newcommand{\resumeProject}[4]{
\vspace{0.5mm}\item
    \begin{tabular*}{0.98\textwidth}[t]{l@{\extracolsep{\fill}}r}
        \textbf{#1} & \textit{\footnotesize{#3}} \\
        \footnotesize{\textit{#2}} & \footnotesize{#4}
    \end{tabular*}
    \vspace{-2.4mm}
    \par
}

\newcommand{\resumeSubItem}[2]{\resumeItem{#1}{#2}\vspace{-4pt}}
\renewcommand{\labelitemi}{$\vcenter{\hbox{\tiny$\bullet$}}$}
\newcommand{\resumeSubHeadingListStart}{\begin{itemize}[leftmargin=*,labelsep=0mm]}
\newcommand{\resumeHeadingSkillStart}{\begin{itemize}[leftmargin=*,itemsep=1.7mm, rightmargin=2ex]}
\newcommand{\resumeItemListStart}{\begin{justify}\begin{itemize}[leftmargin=3ex, rightmargin=2ex, noitemsep,labelsep=1.2mm,itemsep=0mm]\small}
\newcommand{\resumeSubHeadingListEnd}{\end{itemize}\vspace{2mm}}
\newcommand{\resumeHeadingSkillEnd}{\end{itemize}\vspace{-2mm}}
\newcommand{\resumeItemListEnd}{\end{itemize}\end{justify}\vspace{-2mm}}
\newcommand{\cvsection}[1]{%
\vspace{2mm}
\begin{tcolorbox}
    \textbf{\large #1}
\end{tcolorbox}
    \vspace{-4mm}
}
\newcolumntype{L}{>{\raggedright\arraybackslash}X}%
\newcolumntype{R}{>{\raggedleft\arraybackslash}X}%
\newcolumntype{C}{>{\centering\arraybackslash}X}%
%---- End of Packages and Functions ------

%-------------------------------------------
%%%%%%  CV STARTS HERE  %%%%%%%%%%%
%%%%%% DEFINE ELEMENTS HERE %%%%%%%
\newcommand{\name}{Agus Fuad Mudhofar}
\newcommand{\course}{Teknik Komputer}
\newcommand{\phone}{8882111890}
\newcommand{\emaila}{agusfuad090@gmail.com}

\begin{document}
\fontfamily{cmr}\selectfont
%----------HEADING-----------------

{
  \begin{tabularx}{\linewidth}{L r}                                                                                                                                                    \\
    \textbf{\Large \name}                           & {\raisebox{0.0\height}{\footnotesize \faPhone}\ +62-\phone}                                                        \\
    {Institut Teknologi Sepuluh Nopember, Surabaya} & \href{mailto:\emaila}{\raisebox{0.0\height}{\footnotesize \faEnvelope}\ {\emaila}}                                 \\
                                                    & \href{https://maedakenji.github.io/Portofolio/}{\raisebox{0.0\height}{\footnotesize \faGithub}\ {Profil GitHub}}             \\
                                                    & \href{https://www.linkedin.com/in/agus4434/}{\raisebox{0.0\height}{\footnotesize \faLinkedin}\ {Profil LinkedIn}}
  \end{tabularx}
}

%-----------EDUCATION-----------
\section{\textbf{Pendidikan}}
\resumeSubHeadingListStart
\resumeSubheading
{Sarjana Teknik Komputer}{IPK: 3.64}
{Institut Teknologi Sepuluh Nopember, Surabaya}{2022\,--\,2026}
\resumeSubHeadingListEnd
\vspace{-5.5mm}

%-----------EXPERIENCE-----------------

\section{\textbf{Pengalaman}}
\resumeSubHeadingListStart

\resumeSubheading
{Magang Riset -- Pengembang Mobile \& Computer Vision Engineer}{Badan Riset dan Inovasi Nasional (BRIN)}
{Bandung}{Feb -- Juli 2026}

\vspace{0mm}
\begin{itemize}[leftmargin=1.5ex, label={}]

\item{a. Pengembangan Aplikasi Full-Stack}
    \begin{itemize}[label=-]
        \item {Membangun aplikasi Flutter lintas platform (Mobile, Tablet, Desktop) untuk menghitung kebutuhan nutrisi NPK tanaman sayuran, dideploy ke berbagai ukuran layar.}
        \item {Membangun backend RESTful API dengan Laravel 12 dan PostgreSQL; mengimplementasikan autentikasi berbasis token menggunakan Laravel Sanctum.}
        \item {Merancang UI yang intuitif dan responsif di Flutter, memastikan pengalaman pengguna yang mulus di semua perangkat.}
    \end{itemize}

\item{b. Sistem AI \& Computer Vision}
    \begin{itemize}[label=-]
        \item {Mengembangkan sistem visual awareness menggunakan input kamera dan model deep learning untuk mengekstrak informasi kontekstual lingkungan, termasuk atribut pengguna dan objek sekitar.}
        \item {Mengintegrasikan output computer vision dengan Large Language Model (LLM) untuk menghasilkan pemahaman kontekstual yang lebih kaya, memungkinkan interpretasi aktivitas pengguna seperti berjalan, bersepeda, dan berkendara.}
        \item {Merancang modul khusus untuk pipeline pemrosesan yang mengintegrasikan deteksi objek, pengenalan aksi manusia, dan semantic reasoning untuk meningkatkan situational awareness.}
    \end{itemize}
\end{itemize}

\vspace{-2mm}
\resumeSubheading
{Pengembang Web -- Halaman Web MIOT}
{Multimedia and IoT Laboratory (MIOT)}
{Surabaya}{Agu 2024 -- Sekarang}
\resumeItemListStart
\item {Merancang dan mengembangkan halaman web resmi laboratorium MIOT untuk menampilkan penelitian, publikasi, dan anggota tim.}
\item {Mengimplementasikan antarmuka modern dan responsif menggunakan React.js dan Tailwind CSS, memastikan aksesibilitas di berbagai perangkat.}
\item {Mengintegrasikan manajemen konten dinamis untuk berita, acara, dan profil anggota menggunakan struktur data berbasis JSON.}
\item {Mengoptimalkan performa situs dan SEO, serta menyediakan dokumentasi untuk pemeliharaan ke depannya.}
\resumeItemListEnd

\vspace{-2mm}
\resumeSubheading
{Pengembang IoT -- Smart Wheelchair}{Multimedia and IoT Laboratory (MIOT)}
{Surabaya}{Agu 2024 -- Sekarang}
\resumeItemListStart
\item {Memelihara dan meningkatkan sistem kursi roda pintar menggunakan computer vision berbasis deep learning untuk navigasi hands-free.}
\item {Mengintegrasikan pengenalan tatapan mata dan gesture wajah secara real-time menggunakan Python dan OpenCV.}
\item {Mengembangkan firmware mikrokontroler (C/C++) untuk aktuasi motor yang andal dan umpan balik sensor.}
\resumeItemListEnd

\resumeSubheading
{Magang Riset -- Lab Multimedia dan IoT (B300)}{Multimedia and IoT Laboratory (MIOT)}
{Surabaya}{April -- Juni 2025}
\resumeItemListStart
\item {Melakukan riset tentang deployment model deep learning pada edge device menggunakan NVIDIA Jetson Nano, berfokus pada tugas computer vision.}
\item {Mengembangkan dan mengoptimalkan pipeline deteksi objek menggunakan Python dan C++ dengan arsitektur MobileNet-SSD, YOLOv5n, dan YOLOv11s.}
\item {Menjelajahi teknik konversi dan akselerasi model dengan ONNX dan NVIDIA TensorRT untuk inferensi real-time.}
\item {Mengintegrasikan dan mengevaluasi model dalam NVIDIA DeepStream SDK untuk analitik video yang efisien.}
\resumeItemListEnd



\resumeSubheading
{Asisten Pengajar Lab -- Pemrograman Dasar}{Multimedia and IoT Laboratory (MIOT)}
{Surabaya}{Agu 2025 -- Nov 2025}
\resumeItemListStart
\item {Mengembangkan dan memelihara materi lab, termasuk soal latihan dan panduan troubleshooting, sekaligus bertanggung jawab mendeploy dan memelihara website online judge DMOJ untuk perkuliahan.}
\item {Membimbing mahasiswa dalam pengerjaan tugas praktikum C/C++, memberikan dukungan real-time dan penjelasan teknis.}
\item {Mengevaluasi submisi mahasiswa, memberikan umpan balik konstruktif, dan berkontribusi pada rubrik penilaian.}
\resumeItemListEnd
\vspace{-2mm}

% \resumeSubheading
% {Perwakilan -- Demo Kursi Roda Pintar Toyota Exhibition}
% {Toyota Motor Manufacturing Indonesia}
% {Jakarta}{26 -- 27 Mei 2025}
% \resumeItemListStart
% \item {Mewakili proyek kursi roda pintar di Toyota Exhibition, menampilkan inovasi universitas dalam teknologi asistif.}
% \item {Membantu persiapan teknis, termasuk integrasi sistem dan pembaruan perangkat lunak untuk operasi yang andal.}
% \item {Berinteraksi dengan pengunjung, profesional industri, dan perwakilan Toyota, menjelaskan sistem kontrol berbasis computer vision dan dampak sosialnya.}
% \resumeItemListEnd
% \vspace{-4mm}

% End of Experience
\resumeSubHeadingListEnd
\vspace{-5.5mm}

%-----------PROJECTS-----------------

\section{\textbf{Proyek}}
\resumeSubHeadingListStart

\resumeProject
{Aplikasi Web Membaca Komik}
{Platform baca komik berbasis web yang mendukung manga, manhua, dan manhwa dengan pengalaman membaca yang dioptimalkan.}
{}
{}
\resumeItemListStart
\item {Mengembangkan platform komik responsif dengan navigasi halaman yang mulus dan rendering gambar yang dioptimalkan.}
\item {Mengimplementasikan fitur utama termasuk manajemen chapter, riwayat baca, autentikasi, dan layout adaptif untuk berbagai perangkat.}
\item {Teknologi: React (frontend), Laravel (backend/API), PostgreSQL (database), arsitektur komponen modular.}
\resumeItemListEnd
\vspace{-2mm}

\resumeProject
{Sappy: Aplikasi Manajemen Ternak Multi-Peran (Flutter)}
{Aplikasi manajemen ternak untuk berbagai peran (Petani, Admin, Dokter) dengan dukungan NFC.}
{}
{}
\resumeItemListStart
\item {Menyediakan akses berbasis peran dan dashboard untuk petani, administrator, dan dokter dalam satu aplikasi mobile terpadu.}
\item {Mendukung bahasa Inggris dan Indonesia menggunakan lokalisasi bawaan Flutter.}
\item {Mengimplementasikan state management via Provider dan mengaktifkan fungsionalitas NFC untuk identifikasi dan pelacakan aman.}
\item {Teknologi: Flutter, Dart, Provider, NFC, fl\_chart, carousel\_slider.}
\resumeItemListEnd
\vspace{-2mm}

\resumeProject
{Toolkit Kalibrasi Kamera \& Pemrosesan Citra}
{Toolkit untuk kalibrasi kamera, pemrosesan citra, dan eksperimen computer vision menggunakan Python dan React.}
{}
{}
\resumeItemListStart
\item {Mengimplementasikan kalibrasi kamera menggunakan Python dan OpenCV untuk pengukuran gambar yang akurat.}
\item {Mengembangkan Harris Corner Detection dan eksperimen CNN manual untuk klasifikasi gambar.}
\item {Menyediakan frontend React + TypeScript + Vite untuk visualisasi interaktif.}
\item {Teknologi: Python, React, TypeScript, OpenCV, NumPy, Vite.}
\resumeItemListEnd
\vspace{-2mm}

% End of Projects
\resumeSubHeadingListEnd
\vspace{-4mm}


\section{\textbf{Sertifikasi}}
\resumeSubHeadingListStart
\resumeSubheading
{Sertifikasi Kefasihan Bahasa Inggris -- CEFR Level B2}{}
{DynEd International}{2024}
\resumeItemListStart
\item {Telah mencapai kefasihan bahasa Inggris setara dengan Common European Framework of Reference for Languages (CEFR) Level B2.}
\resumeItemListEnd
\resumeSubHeadingListEnd
\vspace{-16pt}


%-----------POSITIONS OF RESPONSIBILITY-----------

\section{\textbf{Tanggung Jawab Organisasi}}
\resumeSubHeadingListStart

\resumeSubheading
{Kompetisi GEMASTIK XVII (Divisi IoT)}
{Universitas Negeri Semarang (UNNES)}{Semarang}
{Jan -- Nov 2024}
\resumeItemListStart
\item {Memimpin tim yang mewakili ITS pada Divisi Smart Devices, Embedded Systems, dan IoT di GEMASTIK XVII.}
\item {Mengembangkan sistem kursi roda pintar yang dapat dikontrol melalui gerakan mata, mengintegrasikan computer vision dan hardware embedded.}
\item {Meraih penghargaan Finalis Nasional (Juara Harapan) melalui perencanaan dan eksekusi inovasi berbasis IoT.}
\resumeItemListEnd
\vspace{-4mm}

\resumeSubheading
{PIC Divisi IoT -- TIM Kawal GEMASTIK}
{Institut Teknologi Sepuluh Nopember (ITS)}{Surabaya}
{Jan 2025 -- Nov 2025}
\resumeItemListStart
\item {Bertugas sebagai Person-in-Charge (PIC) divisi IoT, mengkoordinasikan komunikasi antara pembimbing, dosen, mentor, dan tim peserta.}
\item {Membantu seleksi dan evaluasi tim kandidat melalui review proposal dan asesmen teknis.}
\item {Memfasilitasi sesi mentoring bersama dosen dan pakar domain di bidang IoT.}
\resumeItemListEnd
\vspace{-4mm}

% End of Positions
\resumeSubHeadingListEnd

%-----------SKILLS & INTERESTS-----------------

\section{\textbf{Keahlian Teknis}}
\begin{itemize}[leftmargin=0.05in, label={}]
  \small{\item{
        \textbf{Bahasa Pemrograman}{: C/C++, Python, JavaScript, HTML+CSS, Dart} \\
        \textbf{Framework \& Library}{: React.js, Flutter, Laravel, FastAPI, Tailwind CSS, OpenCV, NumPy} \\
        \textbf{Database}{: PostgreSQL, MySQL, MongoDB, Firebase} \\
        \textbf{Tools}{: Git, GitHub, VS Code, NVIDIA TensorRT, ONNX, Docker (dasar)} \\
        \textbf{AI/ML}{: Deep Learning, Computer Vision, Object Detection (YOLO, MobileNet), LLM Integration} \\
        \textbf{Mata Kuliah Relevan}{: Struktur Data \& Algoritma, Sistem Operasi, OOP, Manajemen Basis Data, Rekayasa Perangkat Lunak} \\
        \textbf{Bahasa}{: Indonesia (Native), Inggris (B2 CEFR - DynEd Certified)} \\
        \textbf{Soft Skills}{: Pemecahan Masalah, Kolaborasi Tim, Komunikasi Teknis, Adaptabilitas, Self-learning} \\
        }}
\end{itemize}


\end{document}
